\documentclass[conference, onecolumn]{IEEEtran}
\usepackage[utf8]{inputenc}
\usepackage[portuguese]{babel}
\usepackage{graphicx}
\usepackage{float}
\usepackage{listings}
\usepackage{xcolor}
\usepackage{booktabs}
\usepackage{amsmath}
\usepackage{caption}
\usepackage{longtable}
\usepackage{pgfplots}
\usepackage{pgfplotstable}
\pgfplotsset{compat=1.18}
\usepackage[hidelinks]{hyperref}

\lstset{
    basicstyle=\ttfamily\footnotesize,
    backgroundcolor=\color{gray!10},
    frame=single,
    breaklines=true,
    columns=fullflexible,
    keepspaces=true,
    tabsize=4,
    showstringspaces=false,
    numbers=left,
    numberstyle=\tiny\color{gray},
    stepnumber=1,
    numbersep=5pt
}

% =========================================================
\title{Trabalho Prático --- Logaritmo Discreto\\
\large Confiabilidade e Segurança de Software (98G08-04)}

\author{
    \IEEEauthorblockN{Bruno Becker Silva}
    \IEEEauthorblockA{Escola Politécnica, PUCRS \\ \texttt{bruno.becker@edu.pucrs.br}}
}

\begin{document}
\maketitle

% =========================================================
\section*{Resumo}

Este relatório documenta o ataque ao Problema do Logaritmo Discreto (DLP) em $({\mathbb Z}/p{\mathbb Z})^*$
para 23 cenários de Diffie--Hellman, com $p$ variando de 16 a 1024 bits.
Foram implementadas três abordagens em C++ com GMP: Força Bruta, Baby-step Giant-step (BSGS)
e Pohlig--Hellman. A chave compartilhada $K_{ab}$ foi recuperada em 15 dos 23 cenários dentro
do limite de 1 hora por cenário.

% =========================================================
\section{Tarefa 1 — Ataques ao DLP: síntese teórica}

\subsection{a) Força Bruta --- $O(p)$}

Testa $x = 0, 1, 2, \ldots$ até encontrar $\alpha^x \equiv A \pmod{p}$.
Trivial de implementar, mas intratável: para $p$ de 32 bits já demora centenas de segundos
num processador moderno; para 64 bits seria astronomicamente mais lento.

\subsection{b) Baby-step Giant-step (BSGS) --- $O(\sqrt{q})$ tempo e memória}

Shanks (1971). Dado $q = \mathrm{ord}(\alpha)$, define $m = \lceil\sqrt{q}\rceil$ e
fatoriza $x = im - j$.
\begin{itemize}
    \item \textbf{Baby steps:} pré-computa a tabela $\{\alpha^j \bmod p : 0 \le j < m\}$.
    \item \textbf{Giant steps:} para $i = 0, 1, \ldots, m$, verifica se
          $A \cdot (\alpha^{-m})^i \bmod p$ está na tabela.
\end{itemize}
Complexidade $O(\sqrt{q})$ em tempo e memória. Para $q \approx p$, resolve o DLP em 64 bits
em segundos, mas exige RAM proporcional a $\sqrt{p}$, tornando-o impraticável acima de
$\sim$80 bits sem estrutura especial em $p-1$.

\subsection{c) Pollard's Rho --- $O(\sqrt{q})$ tempo, $O(1)$ memória}

Usa uma caminhada pseudoaleatória no grupo para detectar colisões via algoritmo
Floyd ou Brent. Mesma complexidade assintótica que BSGS, mas elimina o custo de memória.
É a escolha prática para grupos de ordem prima grande, amplamente usado em ataques reais
a curvas elípticas (ECDLP).

\subsection{d) Pohlig--Hellman}

Reduz o DLP de ordem $q$ ao DLP nos subgrupos de ordem $p_i^{e_i}$ de cada fator primo
de $q = \prod p_i^{e_i}$, e combina via Teorema Chinês do Resto (CRT).
Se $q$ é \emph{smooth} (todos os fatores primos pequenos), a solução é quase instantânea
mesmo para $p$ de 1024 bits. Por isso os parâmetros bem escolhidos de DH exigem que
$p-1$ contenha um fator primo grande (primo seguro ou primo de Sophie Germain).

\subsection{e) Index Calculus / Number Field Sieve (NFS) --- subexponencial}

Complexidade $L_p[\tfrac{1}{3}, c] = e^{(c+o(1))(\ln p)^{1/3}(\ln \ln p)^{2/3}}$.
Utiliza relações entre elementos do grupo e uma base de fatores primos pequenos.
É o algoritmo de referência para quebrar DH em $(\mathbb{Z}/p\mathbb{Z})^*$ com $p$
acima de 1024 bits; quebrou DH-768 em 2016 e DH-1024 está ao alcance de adversários
estatais. Requer infraestrutura de computação distribuída.

% =========================================================
\section{Tarefa 2 — Implementação}

Três soluções foram implementadas em C++17 com aritmética de precisão arbitrária via
\textbf{GMP} (\texttt{mpz\_class}):

\begin{enumerate}
    \item \textbf{\texttt{dlp\_solver\_fb.cpp}} --- Força Bruta: loop simples
          $x = 0, 1, 2, \ldots$, testa $\alpha^x \equiv A \pmod{p}$.
    \item \textbf{\texttt{dlp\_solver\_bsgs.cpp}} --- BSGS: baby steps numa
          \texttt{std::map}, giant steps consultando a tabela.
    \item \textbf{\texttt{dlp\_solver.cpp}} --- Pohlig--Hellman otimizado:
          fatora $p-1$ (Pollard's Rho + crivo), aplica BSGS em cada subgrupo,
          combina via CRT. Multithreaded (até $N_{\mathrm{cpu}}-1$ threads).
\end{enumerate}

O ataque é executado com \texttt{nice -n 10} para ceder prioridade ao sistema operacional.
O código está disponível no repositório do projeto (ver apêndice).

% =========================================================
\section{Tarefa 3 — Resultados e tempos de execução}

A Tabela~\ref{tab:tempos} consolida os tempos medidos das três implementações.
``---'' indica que a execução foi interrompida antes de concluir (timeout $>$ 1h ou
memória insuficiente para BSGS).

\begin{table}[H]
\centering
\caption{Tempos de resolução por cenário e por método. PH = Pohlig--Hellman.}
\label{tab:tempos}
\begin{tabular}{@{}llrrrl@{}}
\toprule
\textbf{Cenário} & \textbf{bits($p$)} & \textbf{FB (s)} & \textbf{BSGS (s)} & \textbf{PH (s)} & \textbf{$K_{ab}$} \\
\midrule
C0  & 16  & 0{,}001   & 0{,}0005  & 0{,}000  & 15261 \\
C1  & 32  & 291{,}0   & 0{,}045   & 0{,}000  & 2315847872 \\
C2  & 32  & 171{,}2   & 0{,}023   & 0{,}029  & 211958971 \\
C3  & 40  & ---       & 0{,}952   & 1{,}781  & 298509300911 \\
C4  & 40  & ---       & 0{,}869   & 0{,}000  & 290240289861 \\
C5  & 48  & ---       & 45{,}10   & 0{,}000  & 87559281402002 \\
C6  & 48  & ---       & 19{,}78   & 42{,}467 & 139741286422486 \\
C7  & 64  & ---       & ---       & 0{,}000  & 5423980969977744708 \\
C8  & 64  & ---       & ---       & $>$3600  & --- \\
C9  & 80  & ---       & ---       & $>$4345  & --- \\
C10 & 80  & ---       & ---       & 0{,}001  & 2551440541766016557167 \\
C11 & 96  & ---       & ---       & 0{,}001  & 27055793202008997473656732406 \\
C12 & 96  & ---       & ---       & $>$4349  & --- \\
C13 & 112 & ---       & ---       & $>$4349  & --- \\
C14 & 112 & ---       & ---       & 0{,}001  & 3555655507609468817009304130465647 \\
C15 & 128 & ---       & ---       & $>$4348  & --- \\
C16 & 128 & ---       & ---       & 0{,}002  & 7831351790774517194506812766211328849 \\
C17 & 256 & ---       & ---       & 0{,}013  & 8631587\ldots855379 \\
C18 & 256 & ---       & ---       & $>$4351  & --- \\
C19 & 512 & ---       & ---       & 0{,}099  & 3567226\ldots774996532 \\
C20 & 512 & ---       & ---       & $>$3600  & --- \\
C21 & 1024& ---       & ---       & $>$3600  & --- \\
C22 & 1024& ---       & ---       & 2{,}670  & 21964133\ldots194578288 \\
\midrule
\multicolumn{5}{l}{\textbf{Total resolvidos:} 15/23 (PH); 7/23 (BSGS); 1/23 (FB)} \\
\bottomrule
\end{tabular}
\end{table}

% =========================================================
\section{Tarefa 4 — Gráfico tempo $\times$ bits($p$)}

\begin{figure}[H]
\centering
\begin{tikzpicture}
\begin{axis}[
    width=\linewidth,
    height=7cm,
    xlabel={bits($p$)},
    ylabel={Tempo (s)},
    ymode=log,
    xmin=10, xmax=1100,
    ymin=1e-4, ymax=1e8,
    xtick={16,32,40,48,64,80,96,112,128,256,512,1024},
    xticklabel style={rotate=45, anchor=east, font=\footnotesize},
    grid=both,
    grid style={line width=0.2pt, draw=gray!30},
    legend style={at={(0.02,0.98)}, anchor=north west, font=\footnotesize},
    legend cell align=left,
]

% --- Força Bruta (medidos) ---
\addplot[
    color=red, mark=*, mark size=2pt, thick,
] coordinates {
    (16,  0.001)
    (32,  231)
};
\addlegendentry{Força Bruta (medido)}

% --- Força Bruta (projeção O(p)) ---
\addplot[
    color=red, dashed, thick, domain=32:80, samples=20,
] {231 * 2^(x-32)};
\addlegendentry{Força Bruta (projeção $O(p)$)}

% --- BSGS (medidos) ---
\addplot[
    color=blue, mark=square*, mark size=2pt, thick,
] coordinates {
    (16,  0.0005)
    (32,  0.034)
    (40,  0.91)
    (48,  32.4)
};
\addlegendentry{BSGS (medido)}

% --- BSGS (projeção O(sqrt(p))) ---
\addplot[
    color=blue, dashed, thick, domain=48:96, samples=20,
] {32.4 * 2^((x-48)/2)};
\addlegendentry{BSGS (projeção $O(\sqrt{p})$)}

% --- Pohlig-Hellman (medidos, cenários fáceis) ---
\addplot[
    color=teal, mark=triangle*, mark size=2.5pt, thick,
] coordinates {
    (16,  0.0001)
    (32,  0.015)
    (40,  0.891)
    (48,  21.23)
    (80,  0.001)
    (96,  0.001)
    (112, 0.001)
    (128, 0.002)
    (256, 0.013)
    (512, 0.099)
    (1024,2.670)
};
\addlegendentry{Pohlig--Hellman (medido)}

% --- Pohlig-Hellman difíceis (timeout) ---
\addplot[
    color=orange, mark=x, mark size=3pt, only marks,
] coordinates {
    (64,  3600)
    (80,  4345)
    (96,  4349)
    (112, 4349)
    (128, 4348)
    (256, 4351)
    (512, 3600)
    (1024,3600)
};
\addlegendentry{Pohlig--Hellman (timeout $>$1h)}

\end{axis}
\end{tikzpicture}
\caption{Tempo de execução (escala log no eixo $y$) por número de bits de $p$.
Os pontos em laranja ($\times$) representam cenários onde o timeout de 1 hora foi atingido.
A dispersão vertical para o mesmo $\log_2 p$ reflete a variação na suavidade (smoothness)
de $p-1$: cenários com $p-1$ liso resolvem em milissegundos; cenários com fator primo
grande em $p-1$ são intratáveis para Pohlig--Hellman.}
\label{fig:tempos}
\end{figure}

% =========================================================
\section{Tarefa 5 — Infraestrutura computacional}

\begin{table}[H]
\centering
\caption{Especificações do ambiente de execução.}
\begin{tabular}{@{}ll@{}}
\toprule
\textbf{Componente} & \textbf{Especificação} \\
\midrule
CPU       & Apple M2 (ARM), 10 núcleos (8P + 2E) \\
Threads   & Até 9 threads simultâneas (1 núcleo reservado ao SO) \\
RAM       & 16 GB LPDDR5 unificada \\
SO        & macOS 15 (Sequoia), kernel Darwin \\
Compilador& \texttt{g++ -std=c++17 -O3 -march=native} \\
Biblioteca& GMP 6.x (\texttt{mpz\_class}, precisão arbitrária) \\
Prioridade& \texttt{nice -n 10} (baixa prioridade de escalonamento) \\
Ataque    & Local, single-node (sem distribuição) \\
\bottomrule
\end{tabular}
\end{table}

% =========================================================
\section{Tarefa 6 — Discussão crítica}

\subsection{O que os resultados revelam sobre escolha de parâmetros em DH real}

\paragraph{A ordem do subgrupo, não o tamanho de $p$, determina a dificuldade.}
Os cenários C7 (64 bits), C10, C11, C14, C16, C17, C19 e C22 foram resolvidos em menos
de 3 segundos com Pohlig--Hellman, enquanto C8 (também 64 bits), C9, C12, C13, C15, C18
e C20 ultrapassaram 1 hora. A diferença está na fatoração de $p-1$: nos cenários fáceis,
$p-1$ é \emph{liso} (todos os fatores primos pequenos), reduzindo o DLP a subproblemas
triviais. Nos difíceis, $p-1$ contém um fator primo grande, tornando Pohlig--Hellman
equivalente a resolver o DLP na ordem total.

\paragraph{BSGS é prático até $\sim$50 bits, depois a memória é o gargalo.}
O BSGS resolveu C0--C6 (até 48 bits) em tempo razoável, mas falhou em C7 (64 bits) por
exigir $2^{32}$ entradas na tabela hash ($\sim$64 GB de RAM). Isso ilustra que algoritmos
de memória proporcional a $\sqrt{p}$ são impraticáveis acima de 60--70 bits em hardware
convencional.

\paragraph{Força Bruta é apenas pedagógica.}
O método resolveu C0 (16 bits) trivialmente e C1/C2 (32 bits) em 3--5 minutos. Para
C3 (40 bits) o tempo estimado seria $\sim$290 horas; para 64 bits, superior à idade
do universo.

\paragraph{Parâmetros seguros em DH real.}
NIST e RFC 3526/7919 especificam grupos com $p \geq 2048$ bits e $p-1 = 2q$ com $q$
primo (\emph{primo seguro}), garantindo que a ordem máxima do subgrupo seja $q \approx p/2$.
Nesse regime, mesmo Pollard's Rho precisaria de $\sim 2^{1024}$ operações, inatingível
computacionalmente. Os experimentos confirmam que parâmetros mal escolhidos (p-1 liso)
colapsam a segurança mesmo para $p$ de 1024 bits.

\paragraph{Index Calculus e NFS: o horizonte real.}
Para DH em $(\mathbb{Z}/p\mathbb{Z})^*$ com $p$ grande e $p-1$ com fator primo grande,
o ataque viável é Index Calculus / NFS. A complexidade subexponencial
$L_p[1/3, c]$ torna 1024 bits marginalmente seguro contra adversários com grande
capacidade computacional (DH-1024 foi considerado quebrado em cenários de computação
distribuída de longo prazo). A migração para grupos de curvas elípticas (ECDH) com 256
bits oferece segurança equivalente a DH com 3072 bits em $(\mathbb{Z}/p\mathbb{Z})^*$.

% =========================================================
\section*{Conclusão}

Os experimentos demonstraram empiricamente que a segurança do protocolo Diffie--Hellman
\emph{não depende apenas do tamanho de $p$}, mas fundamentalmente da estrutura de $p-1$.
A implementação de Pohlig--Hellman resolveu 15 dos 23 cenários, incluindo exemplos de
1024 bits com $p-1$ liso, em menos de 3 segundos. Isso torna concreto o requisito
criptográfico de utilizar primos seguros ou grupos padronizados com fator primo grande
em qualquer aplicação séria de Diffie--Hellman.

% =========================================================
\appendix
\section*{Apêndice — Código-fonte}

O código está organizado em três arquivos principais:

\begin{itemize}
    \item \texttt{dlp\_solver\_fb.cpp} --- Força Bruta com salvamento incremental e
          tratamento de \texttt{SIGINT}.
    \item \texttt{dlp\_solver\_bsgs.cpp} --- Baby-step Giant-step com \texttt{std::map}
          e salvamento incremental.
    \item \texttt{dlp\_solver.cpp} --- Pohlig--Hellman com Pollard's Rho para fatoração,
          BSGS nos subgrupos, CRT e multithreading.
\end{itemize}

Compilação (macOS com GMP via Homebrew):
\begin{lstlisting}[language=bash, numbers=none]
make all          # compila todos os tres executaveis
make runfb        # roda forca bruta em desafios.txt
make runbsgs      # roda BSGS em desafios.txt
make run          # roda Pohlig-Hellman em desafios.txt
\end{lstlisting}

\end{document}
